% October 30, 2007
% Math 589

\documentclass[
    10pt % font size
    16:9, % 1920x1080
]{beamer}
% \usetheme{default}
% \usetheme{Boadilla}
 \usetheme{Madrid}
% \usetheme{Montpellier}
% \usetheme{Warsaw}
% \usetheme{Copenhagen}
% \usetheme{Goettingen}
% \usetheme{Hannover}
% \usetheme{Berkeley}
 
% \usecolortheme{crane}
 % \beamertemplatesolidbackgroundcolor{craneorange!25}
 
 % Define custom colors
\definecolor{customGreen}{RGB}{0,128,0} % A medium green
\definecolor{customDarkGreen}{RGB}{0,100,0} % A darker green

\usepackage{hyperref}

% Apply the custom colors
\usecolortheme{default} % Start with the default color theme to apply custom colors on top
\setbeamercolor{structure}{fg=customGreen}
\setbeamercolor{background canvas}{bg=white} % Set main background to white
\setbeamercolor{title}{bg=customDarkGreen,fg=white} % Title slide background
\setbeamercolor{frametitle}{bg=customGreen,fg=white} % Frame titles

% Custom footline with an image on the bottom right
\addtobeamertemplate{footline}{}{%
  \hfill%
  \raisebox{5mm}[0pt][10pt]{%
    \includegraphics[height=1cm]{kyu_univ_logo.png}%
  }\hspace*{5mm}
}

\title{Article Classification with Centroid Classification}
\subtitle{}

\author{Group 1}

\date[June 19, 2024]
 {June 19, 2024}

\begin{document}

\frame{\titlepage} % # 1
\section[Outline]{}
\frame{\tableofcontents}

\section{Project Overview}
 
\frame % # 2
{
  \frametitle{Project Objectives}
  \begin{itemize}
    \item The goal of this project is classify document per category on given dataset.
    \item 
  \end{itemize}
}
\section{Dataset and Pre-liminary Analysis}
\frame
{
  \frametitle{Dataset and Pre-liminary Analysis}
  \begin{itemize}
    \item \href{http://mlg.ucd.ie/datasets/bbc.html}{BBC ML Dataset}
    \item Consist of 62225 articles, with categories: 'business', 'entertainment', 'politics', 'sport', 'tech'
  \end{itemize}
}

\section{Data Preprocessing}
\frame
{
  \frametitle{Data Preprocessing}
  \begin{itemize}
    \item Data is loaded from .parquet files
    \item Data is cleaned by removing special characters, numbers, and stopwords
    \item Stopwards are set of words that are considered unimportant in the text data. For stop word, we used nltk package.
    \item After cleaning, data is converted into lower-case where it's used to generate a unique words list.
    \item Documents are converted into numerical data using Bag of Words method
  \end{itemize}
}

\begin{frame}
  \frametitle{Bag of Words}
  \begin{itemize}
    \item The Bag of Words method converts text data into numerical data by representing each document as a vector of word counts.
  \end{itemize}
  \[
    \mathbf{d} = \begin{bmatrix}
    f_1 \\
    f_2 \\
    \vdots \\
    f_n
  \end{bmatrix}
  \]
  \begin{itemize}
    \item Here, \(\mathbf{d}\) is a vector representing a document, and \(f_i\) is the frequency of the \(i\)-th word from the vocabulary in the document.
  \end{itemize}
\end{frame}

\begin{frame}
  \frametitle{tf-idf scoring and minimum frequency}
  \begin{itemize}
    \item After BoW Encoding, we calculate tf-idf score for each token and documents.
    \item We drop tokens with frequency less than 20, and tf-idf score less than 30.
    \item For tf-idf, we use the following formula:
    \begin{align}
      tf(t,d) &= 0.5 + 0.5 \times \frac{f_{t,d}}{\max_{t' \in d} f_{t',d}} \\
      idf(t) &= \log\left(\frac{N}{n_t}\right) \\
      tfidf(t,d) &= tf(t,d) \times idf(t)
    \end{align}
  \end{itemize}
\end{frame}

\section{Document Vectorization and Clustering}
\begin{frame}
  \frametitle{Document Vectorization and Clustering}
  \begin{itemize}
    \item After all documents are vectorized, we apply dimension reduction(CUR decomposition) to reduce the dimension of the data.
    \item CUR decomposition is utilized by sparse nature of embedded vector.
    \item Dataset is batched in to 1500 rows per batch, and CUR decomposition is applied to each batch.
    \item Afterwards, per unique category we may construct a probability distribution using the centroid of the cluster and distance metric.
  \end{itemize}
\end{frame}

\begin{frame}
  \frametitle{Document Cluster and probability distribution}
  \begin{itemize}
    \item For each unique category, we construct a probability distribution using the centroid of the cluster and distance metric.
    \item Assuming central-limit theorem, we use a normal distribution to generate the probability distribution.
    \[
      p(d) = \frac{1}{\sqrt{2\pi}\sigma}e^{-\frac{(d)^2}{2\sigma^2}}
    \]
    Where d is the distance from the centroid, and \(\sigma\) is the variance of the cluster.
  \end{itemize}
\end{frame}

\begin{frame}
  \frametitle{Pmf Training}
  \begin{itemize}
    \item After generating the probability distribution, we use it to predict the tag of a given wikipedia document.
    \item We can tune the probability threshhold for accuracy using the cost function:
    \[
      J(C_d, C_d') = \frac{\lambda_1|C_d'\setminus C_d| + \lambda_2|C_d\setminus C_d'|}{|C_d\cup C_d'|}
    \]
    where $C_d$ is document categories set, and $C_d'$ is predicted categories set.
    \item Starting from threshold 0, we increase the threshold by some small value until the cost function is minimized.
  \end{itemize}
\end{frame}

\section{Current Progress}
\begin{frame}
  \frametitle{Current Progress}
  \begin{itemize}
    \item Currently generated tf-idf score for all token using 10\% sampled data, and filteered out out tokens with low score.
    \item CUR decomposition is in process of implementation, and will soon be applied to all data.
    \item After CUR decomposition, we need to split the data into training and testing sets, and construct a probability distribution for each unique category.
    \item For efficiency, we plan to implement further computation with rapids library for GPU acceleration, using cuml and cudf, and pytorch.
  \end{itemize}
\end{frame}
\end{document}